%%%%%%%%%%%%%%%%%%%%%%%%%%%%%%%%%%%%%%%%%%%%%%%%%%%%%%
%%
%%           EE474 C constants and #define
%%
%%
%   % combined into notes pack
%
%\documentclass{article}
%\documentclass{seminar}
%
%\newpagestyle{BHstyle}{\sl U. of Washington \hfil EE 474 \hfil \thepage}{Hannaford \hfil \today}
%\pagestyle{headings}
%\slidepagestyle{BHstyle}
%\addtolength\oddsidemargin{3cm}
%\usepackage[T1]{fontenc}
%\usepackage[utf8]{inputenc}
%\usepackage[english]{babel}
%\usepackage{graphicx}
%\usepackage[pdftex,
%    pdftitle={ECE 474: C Constants and Preprocessor Directives},
%    pdfauthor={Blake Hannaford, University of Washington},
%    pdflang={en-US},
%    pdfsubject={ECE 474 Embedded Microcomputer Systems},
%    pdfkeywords={embedded systems, microcontrollers, ECE 474},
%    hidelinks,
%    unicode]{hyperref}
%\usepackage{listings}
%
%
%\oddsidemargin 0in
%\evensidemargin 0in
%\textwidth 6.5in
%\textheight 8.7in
%\topmargin 0.2in
%
%
%
%
%\linespread{1.0}
%
%
%\begin{document}
	%<*>
%\begin{slide}
\section*{{\tt C} Constants, Defines and Integers \footnote{
 B. Kernigan and D. Ritchie, ``The C Programming Language,,''
 2nd Edition, Prentice Hall, 1988.  \\
 Daniel W. Lewis, Fundamentals of
 Embedded Software, {\it Where C and Assembly Meet}, Prentice Hall, 2002.}
 }

%\end{slide}
%\begin{slide}

%%%%** Section 1 
\section{Review of C Constants}

We will often need to specify the content of memory
very precisely.

%%%%** Section 1.1
\subsection{HEX and OCTAL review}

	%<*chn>
Entering binary numbers can wear out the 1 and 0 keys on
your computer!   HEX is a good shorthand for binary since
each HEX character maps to a specific bit pattern.


\vspace{0.2in}	%<chn>
	%<*>
\begin{tabular}{|c|c|r|}
\hline
bits  &  HEX & Octal \\ \hline
0000  &  0   & 0     \\ \hline
0001  &  1   & 1     \\ \hline
0010  &  2   & 2     \\ \hline
...   &  ...    & ...      \\ \hline
0111  &  7   & 7     \\ \hline
1000  &  8   & 10     \\ \hline
1001  &  9   & 11     \\ \hline
1010  &  A   & 12     \\ \hline
1011  &  B   & 13     \\ \hline

...   & ...  & ...   \\ \hline
1111  &  F   & 17     \\
\hline
\end{tabular}

\vspace{0.2in}	%<chn>

Hex Place value: To convert {\tt 0xD2FB} to decimal:

%\vspace{0.2in}
\begin{tabular}{|c|c|c|c|}
\hline
D       &       2       &       F       &       B \\ \hline
$2^{12}$        &       $2^8$   &       $2^4$   &       $2^0$ \\ \hline
4096    &       256     &       16      &       1 \\ \hline
        &               &               &          \\ \hline
\end{tabular}
%\vspace{0.2in}

	%<*chn>
Octal
can also be used and is easier to memorize because there are
only 8 symbols (0-7). However, to fill
a byte of memory requires 2.5 octal characters since each
character only encodes 3 bits.  For example,
	%<*>

{\tt HEX 9A}  = Binary {\tt 1011010 } = Octal {\tt 132}

(Hex: 4-bit groups.   Octal: 3-bit groups)

%\end{slide}
%\begin{slide}


\begin{itemize}
  \item
  Ordinary numbers are interpreted by the C compiler
          as decimal values.   Example: {\tt 25 = 1101}

  \item Hex values are indicated to the C compiler by the
              prefix {\tt Ox} for example:   {\tt Ox2F}

  \item Octal values are indicated by a leading zero: {\tt 028}
            is a syntax error (why?).

\end{itemize}


%\end{slide}
%\begin{slide}

%%%%** Section 1.2
\subsection{Constant Syntax in C}

It is very important to use symbolic constants ({\tt \#define}'s )
for {\em all} your numeric constants.  Here are the
main reasons:
\begin{itemize}
        \item Code is more readable by humans.
        \item Code is much easier to change.
\end{itemize}

%%%%** Section 1.2.1 
\subsubsection{Example and Exercise}
	%<*chs>
{\small\tt
\begin{verbatim}
// Constants
#define EXAMPLE_CONST_D  1234 ;  /* DECIMAL value  */
#define EXAMPLE_CONST_H 0x4D2 ;  /* HEX value  */
#define EXAMPLE_CONST_O 02322  ; /* OCTAL value */

int x = EXAMPLE_CONST_D;
int y = EXAMPLE_CONST_H;
int z = EXAMPLE_CONST_O;

// test yourself here: (which of these print?)
if (x == y) printf ("Hello there ... \n");
else {};

if (y == 1234) printf ("What's up doc?? \n");
else {};

if (z == 2322)  printf ("The Rain in Spain  ... \n");
else {};

if (z = 2322)  printf ("Falls mainly on the plains.\n");
else {};

\end{verbatim}
}

%
%{\small\tt
%\begin{verbatim}
%// Constants
%#define EXAMPLE_CONST_D  1234 ;  /* DECIMAL value  */
%#define EXAMPLE_CONST_H 0x4D2 ;  /* HEX value == dec: 1234 */
%#define EXAMPLE_CONST_O 02322  ; /* OCTAL value */
%
%int x = EXAMPLE_CONST_D;
%int y = EXAMPLE_CONST_H;
%int z = EXAMPLE_CONST_O;
%
%// test yourself here: (which of these print?)
%if (x == y) printf ("Hello there ... \n");  **YES**
%else {};
%
%if (y == 1234) printf ("What's up doc?? \n");  **YES**
%else {};
%
%if (z == 2322)  printf ("The Rain in Spain  ...\n");  **NO**
%else {};
%
%if (z = 2322)  printf ("Falls mainly on the plains.\n");   **YES**
%else {};
%
%\end{verbatim}
%}
	%<*>


%\end{slide}
%\begin{slide}


%%%%** Section 1.2.2 
\subsubsection{Example 2}
Suppose you are going to control a CD-ROM drive and you want
to open and close the drive door.  The bit to open and close
the drawer is the 4th bit in a register located at address
{\tt 0x2DA}.  Let's say that to set that bit we use the function
\begin{verbatim}
IO_Register_Set([addr],[value])
\end{verbatim}
where {\tt [addr]}
is the bus address of the I/O register you want to manipulate
and {\tt [value]} is the bit pattern you want to put there.

It is ``technically'' correct to use the following code:
\begin{verbatim}
\\ open CD-ROM door
IO_Register_Set(0x2DA, 0x10) ;
\end{verbatim}


%\end{slide}
%\begin{slide}

	%<*cnh>
However, if we come back to that code later (1 week or 10 years)
we will not know what that statement does (or if the comment is
correct) without a lot of research.  Instead we use defines
for much better code as follows:
%Use {\tt\#define} compiler directives instead!
	%<*>
\begin{verbatim}

\\ (this part is in the preamble before "main ()"
\\  or in a special ".h" include file
#define  CD_ROM_CONTROL_REG  0x2DA
#define  CD_ROM_DOOR_OPEN    0x010

\\ .... skip to inside the code

\\ open CD-ROM door
IO_Register_Set(CD_ROM_CONTROL_REG, CD_ROM_DOOR_OPEN) ;

\end{verbatim}

Now, isn't that better?!!
%
It is a requirement of EE474 that
your code be written this way.
%
We will take off points for
any numberical constants in the main body of code.

%\end{slide}
%\begin{slide}

%%%%** Section 2 
\section{Bitwise vs. Logical Operators}

The familiar Boolean
operations {\tt\bf AND}, {\tt\bf OR}  and  {\tt\bf NOT}
are provided {\em two ways} in C.


%%%%** Section 2.1
\subsection{Logical Operators}

These work on {\em arithmetic } types or pointers.

\begin{tabular}{l p{1.5in} }

{\tt\bf True}           & a non zero value. \\

{\tt\bf False}          & zero.           \\

{\tt\bf \&\&}           & {\tt\bf AND}     \\

{\tt A \&\& B }         & True if {\it both} {\tt A} and {\tt B} are non-zero. \\

{\tt ||}                &  {\tt\bf OR} \\

{\tt A || B }           &  True if {\it either} {\tt A} and {\tt B} are non-zero.\\

{\tt ! }                &   {\tt\bf NOT} \\

\end{tabular}


%\end{slide}
%\begin{slide}

%%%%** Section 2.2
\subsection{Bitwise Operators}

A very important C feature for embedded microcomputer systems.
These operators are only defined for integer-like variables i.e.
{\tt char, short, int, and long} signed or unsigned.


\begin{tabular}{l l}
\&                      &       bitwise AND \\
\textbar                &       bitwise OR   \\
\textasciicircum        &       bitwise XOR  \\
\textless\textless      &       left shift   \\
\textgreater\textgreater&       right shift    \\
\textasciitilde         &       ones complement \\

\end{tabular}


%\end{slide}
%\begin{slide}

%%%%** Section 2.3
\subsection{Examples:}

{\tt
\begin{verbatim}

/* Use of bitwise logical operators  */

#define  Bit_Zero   0x01
#define  Bit_One    0x02
#define  Bit_Two    0x04
#define  Bit_Three  0x08
#define  Bit_Four   0x10
#define  Bit_Five   0x20
#define  Bit_Six    0x40

// etc etc etc ..


int x = 0x0B;           // x  = [... 0 0 0 0 1 0 1 1]

int y = 011;            // y  = [... 0 0 0 0 1 0 0 1]

int z = x << 2;         // z  = [... 0 0 1 0 1 1 0 0]

int z1 = y & Bit_Three; // z1 = [... 0 0 0 0 1 0 0 0]

int z2 = z | Bit_Four;  // z2 = [... 0 0 1 1 1 1 0 0]

\end{verbatim}
}


%\end{slide}
%\begin{slide}

\newpage	%<!s>
%%%%** Section 3 
\section{In-class Exercise:}

%
%\begin{verbatim}
%
%Convert 0x1A to decimal: __26_
%
%
%Convert 0xA2 to decimal:  __162____
%
%
%Convert 020 to Hex:    _____10_______
%
%
%Convert 0x20 to Octal: _____40_________
%
%\end{verbatim}
%
%{\tt
%\begin{verbatim}
%/* Give the binary value of the least significant
%8 bits of the following values:   */
%
%int aa = 0xC5;     // aa = [  1 1 0 0 0 1 0 1  ]
%
%int ab = 017 + 1;  // ab = [  0 0 0 1 0 0 0 0  ]
%
%int a = x | y;     // a  = [  0 0 0 0 1 0 1 1  ]
%
%int b = Bit_Two | Bit_Five | Bit_Seven;
%                   // b  = [  1 0 1 0 0 1 0 0  ]
%
%int c = ~((z | Bit_Zero) << 1);
%                   // c  = [  1 0 1 0 0 1 0 1 ]
%
%int d = !(z1 | Bit_Two);
%                   // d  = [  0 0 0 0 0 0 0 0 ]
%\end{verbatim}
%}
%
	%<*chs>

\begin{verbatim}

Convert 0x1A to decimal: __________


Convert 0xA2 to decimal:  __________


Convert 020 to Hex:  _____________________


Convert 0x20 to Octal: _____________________

\end{verbatim}

{\tt
\begin{verbatim}
/* Give the binary value of the least significant
8 bits of the following values:   */

int aa = 0xC5;     // aa = [                     ]

int ab = 017 + 1;  // ab = [                     ]

int a = x | y;     // a  = [                     ]

int b = Bit_Two | Bit_Five | Bit_Seven;
                   // b  = [                     ]

int c = ~((z | Bit_Zero) << 1);
                   // c  = [                     ]

int d = !(z1 | Bit_Two);
                   // d  = [                     ]
\end{verbatim}
}

	%<*>
\newpage	%<!s>
%%%%** Section 4 
\section{Binary Integers}

%%%%** Section 4.1
\subsection{Unsigned Integer}

Typically just a
16 bit binary number.  Straightforward, but cannot represent
negative numbers.

If {\tt x } is of type {\tt unsigned int}, then
\[
0 \leq {\mathtt x} \leq 65535
\]


%\end{slide}
%\begin{slide}

C declaration of unsigned integers:

\begin{verbatim}
unsigned char a,b;            // 8 bit unsigned
unsigned int  x,y,z;          // 16 bit unsigned
unsigned long int p,q,r;      // 32 bit unsigned
\end{verbatim}



%\end{slide}
%\begin{slide}

%%%%** Section 4.1.1 
\subsubsection{Exercise}

Write a C function using bitwise logical operators etc. to print
the binary value of a 16 bit word.

\vspace{4in}	%<cn>

%\end{slide}
%\begin{slide}

%\vspace{3.0in}
%
%\begin{verbatim}
%void bin_display(unsigned int x)
%{
%        unsigned int i,j;
%        for(i=0;i<16;i++)
%        {
%                j = 1<<(15-i);
%                if(j & x) printf("1");
%                else      printf("0");
%        }
%}
%
%\end{verbatim}
%
%\end{slide}
%\begin{slide}
%
	%<*>
%%%%** Section 4.2
\subsection{2's Complement Integers}
	%<*chn>
Most common way to represent integers which can have both positive and negative
values.   Most significant bit (MSB) is the ``sign bit'' 0=positive, 1=negative.
however, if MSB is 1, other bits are subject to ``2's complement''.
%
%\begin{itemize}
%        \item MSB is ``sign bit'' 0= positive, 1=negative
%        \item if sign bit is 1, rest of word is in ``2's complement''
%\end{itemize}
%
	%<*>

To convert a number to or from two's complement,
\begin{enumerate}
  \item complement all bits
  \item add 1
\end{enumerate}
Taking the 2's complement of the entire $n$ bits is equivalent to negation.

For an $n$ bit number, $x$,  the range of all possible integers in 2's complement is
\[
-2^{n-1} \leq x \leq (2^{n-1}-1)
\]
(see Lewis Figure 2-4, page 22 for a useful chart).

%\end{slide}
%\begin{slide}

%%%%** Section 4.2.1 
\subsubsection{Exercise}

Write a c code fragment which negates a signed (2's complement) integer
using bitwise logical operators and addition.  Include declarations for
your variables.
(do not just multiply by -1)


\vspace{3.0in}	%<hc>

%\end{slide}
%\begin{slide}
%
%\begin{verbatim}
%
%[blake@alec code]$ cat twocomp.c
%/* twocomp.c
% *
% *  Demonstrate 2's complement
% *
% */
%void bin_display(unsigned int x);
%main(argc,argv)
%int argc;
%char *argv;
%{
%        int  a = -23;
%        int  b;
%        int i,j;
%        printf(" Decimal: %d \n    Binary:     ", a);
%        b = a;
%        bin_display((unsigned int)b);
%        printf("\n    Complement: ");
%        b =   ~a;     // complement all bits
%        bin_display((unsigned int)b);
%        printf("\n1 + Complement: ");
%        b = 1 + ~a;     // complement all bits and add 1
%        bin_display((unsigned int)b);
%        printf("\n Compare: %d %d\n", a, -b);   // should be the same
%}
%
%
%\end{verbatim}
%\end{slide}
%\begin{slide}
%%%%%** Section 4.3
%\subsection{Test Run}
%\begin{verbatim}
%
%[blake@alec code]$ ./twocomp
% Decimal: -23
%    Binary:     1111111111101001
%    Complement: 0000000000010110
%1 + Complement: 0000000000010111
% Compare: -23 -23
%[blake@alec code]$
%
%\end{verbatim}
%
%\end{slide}
%\end{document}
%
